\chapter{Recomendaciones}

\section{Recomendaciones metodológicas}

Se recomienda evaluar la sensibilidad de los resultados a diferentes métodos de detrending, definiciones de anomalías negativas de rendimiento, índices ENSO, métricas MJO y fuentes de calendarios agrícolas. También se recomienda comparar enfoques basados en reglas interpretables con métodos de clasificación estadística o aprendizaje automático, manteniendo la interpretabilidad como criterio central para aplicaciones en seguros paramétricos.

\section{Recomendaciones para validación futura}

La validación futura debería incluir contraste con datos subnacionales, información de pérdidas reportadas, datos de seguros agrícolas cuando estén disponibles y validación experta por región. También se recomienda evaluar la estabilidad temporal de los umbrales mediante esquemas de validación fuera de muestra.

\section{Recomendaciones para desarrollo de la herramienta web}

La herramienta web debería permitir consultas por cultivo, región, fase fenológica, régimen agroclimático e índice climático. Se recomienda incluir mapas interactivos, descarga de tablas, visualización de series temporales, documentación metodológica y advertencias explícitas sobre incertidumbre y uso responsable de los triggers.

% TODO: definir stack tecnológico sugerido, formato de datos y flujo de actualización.

\section{Recomendaciones para aseguradoras, bancos, gobiernos y entidades de seguridad alimentaria}

Para aseguradoras y bancos, se recomienda usar los resultados como insumo inicial de segmentación del riesgo, no como producto actuarial final. Para gobiernos y entidades de seguridad alimentaria, se recomienda integrar los mapas de sensibilidad y regímenes agroclimáticos en sistemas de alerta temprana, planificación de contingencias y priorización de intervenciones. Para todos los actores, se recomienda evaluar el riesgo base, la transparencia del índice y la pertinencia territorial antes de implementar cualquier mecanismo paramétrico.

\section{Recomendaciones para trabajo doctoral o publicaciones futuras}

Como trabajo futuro, se recomienda ampliar la investigación hacia modelos causales, análisis de eventos compuestos, integración de pronósticos estacionales y subseasonales, evaluación de cambio climático, y diseño formal de productos paramétricos con validación actuarial. También se recomienda preparar publicaciones derivadas sobre regímenes agroclimáticos globales, sensibilidad fenológica a ENSO/MJO y diseño de triggers climáticos para cereales.
